% Part: computability
% Chapter: computability-theory
% Section: coding-computations

\documentclass[../../../include/open-logic-section]{subfiles}

\begin{document}

\olfileid[fr]{cmp}{thy}{cod}
\olsection{Codage des calculs}

Dans les modèles de calcul usuels, on peut accomplir
les opérations suivantes :
\begin{enumerate}
\item Décrire systématiquement les \emph{définitions}
  des fonctions calculables. Les descriptions de machines
  de Turing, les définitions récursives et les programmes
  écrits dans un langage de programmation en sont des exemples.
\item Décrire l'enregistrement complet du calcul d'une
  fonction, donnée par une définition, sur une entrée
  donnée. Pour une machine de Turing, il peut s'agir de
  la suite des configurations à chaque étape : état
  de la machine et contenu du ruban.
\item Vérifier si un enregistrement proposé décrit
  effectivement le calcul d'une fonction ayant une
  définition donnée sur une entrée donnée, pour laquelle
  la fonction est définie, c'est-à-dire où le calcul s'arrête.
\item Extraire d'un enregistrement complet la valeur de
  la fonction sur l'entrée considérée. Pour une machine
  de Turing qui s'arrête, le contenu du ruban à la
  dernière étape fournit, par exemple, cette valeur.
\end{enumerate}

Grâce au codage, on peut attribuer à chaque description
d'une fonction calculable un \emph{indice} numérique
permettant de retrouver effectivement les instructions
qu'elle contient. On peut de même coder par un seul
entier l'enregistrement complet d'un calcul. La relation
arithmétique « $s$ code l'enregistrement d'un calcul terminé de
la fonction d'indice~$e$ sur l'entrée~$x$ » et la
fonction « résultat de la suite de calcul de code~$s$ »
sont alors calculables; en fait, elles sont récursives
primitives.

Ce fait fondamental permet de démontrer plusieurs résultats
importants et parfois surprenants sur la calculabilité,
indépendamment du modèle de calcul choisi parmi ceux
qui admettent un tel codage.\footnote{Précision éditoriale : « every model » dans la source est compris ici comme un modèle effectif usuel muni d'un codage approprié. La relation de vérification porte sur un calcul terminé, pas sur un simple préfixe.}

\end{document}
